\documentclass[10pt, aspectratio=169]{beamer}
\usepackage{../beamer_theme_408}

\title{408考研零基础 --- 操作系统}
\subtitle{3-16 SPOOLing技术}
\author{}
\date{}

\begin{document}

\frame{\titlepage}

\begin{frame}{本节目标}
    \begin{enumerate}
        \item 理解 SPOOLing 技术的核心概念与思想
        \item 掌握 SPOOLing 系统的组成结构
        \item 理解打印机 SPOOLing 的完整工作流程
        \item 理解虚拟设备技术的含义
    \end{enumerate}
\end{frame}

\begin{frame}{SPOOLing 的基本概念}
    \textbf{SPOOLing} = Simultaneous Peripheral Operation On-Line
    \vspace{0.5em}
    \begin{itemize}
        \item 中文：\textbf{假脱机操作}
        \item 在多道程序环境下，利用磁盘作为缓冲
        \item 将\textbf{独占设备}改造为\textbf{共享设备}
    \end{itemize}
    \vspace{0.5em}
    核心思想：
    \[
    \boxed{\text{用磁盘的高速缓冲区模拟低速独占设备}}
    \]
    \vspace{0.5em}
    \begin{itemize}
        \item 用户进程不需要直接操作物理设备
        \item 所有 I/O 请求先写入磁盘，再由 OS 统一调度
    \end{itemize}
\end{frame}

\begin{frame}{SPOOLing 系统的组成}
    SPOOLing 系统由四部分组成：
    \vspace{0.5em}
    \begin{center}
        \begin{tabular}{|l|l|}
            \hline
            \textbf{组件} & \textbf{功能} \\
            \hline
            输入井 & 磁盘区域，暂存输入数据 \\
            输出井 & 磁盘区域，暂存输出数据 \\
            输入进程 & 负责从设备读入数据到输入井 \\
            输出进程 & 负责从输出井写出数据到设备 \\
            输入缓冲区 & 内存缓冲区，缓和设备与磁盘速度差 \\
            输出缓冲区 & 内存缓冲区，缓和磁盘与设备速度差 \\
            \hline
        \end{tabular}
    \end{center}
    \vspace{0.5em}
    输入井和输出井是磁盘上的\textbf{专用存储区域}。
\end{frame}

\begin{frame}{SPOOLing 系统结构图}
    \[
    \begin{array}{c}
        \text{用户进程 1} \quad \text{用户进程 2} \quad \text{用户进程 3} \\
        \downarrow \qquad \downarrow \qquad \downarrow \\
        \boxed{\text{输入缓冲区}} \qquad \boxed{\text{输出缓冲区}} \\
        \downarrow \qquad \downarrow \\
        \boxed{\text{输入井}} \qquad \boxed{\text{输出井}} \\
        \downarrow \qquad \downarrow \\
        \boxed{\text{输入进程}} \qquad \boxed{\text{输出进程}} \\
        \downarrow \qquad \downarrow \\
        \text{输入设备} \qquad \text{输出设备}
    \end{array}
    \]
    \vspace{0.5em}
    用户进程通过 SPOOLing 系统间接访问设备，无需直接操作硬件。
\end{frame}

\begin{frame}{打印机 SPOOLing 案例}
    打印机是典型的\textbf{独占设备}：
    \begin{itemize}
        \item 同一时刻只能由一个进程使用
        \item 多进程直接打印会导致输出混杂
    \end{itemize}
    \vspace{0.5em}
    SPOOLing 解决方案：
    \vspace{0.5em}
    \begin{enumerate}
        \item 用户进程调用 \texttt{print()}，系统将数据写入\textbf{输出井}
        \item 系统为该打印请求创建\textbf{打印请求表}并挂入队列
        \item 用户进程立即返回，感觉"打印完成"
        \item \textbf{打印机守护进程}从队列中取出请求
        \item 从输出井读取数据，送往实际打印机
    \end{enumerate}
    \vspace{0.5em}
    结果：用户进程互不阻塞，打印机被高效共享。
\end{frame}

\begin{frame}{打印机 SPOOLing 示意图}
    传统方式（无 SPOOLing）：
    \[
    \text{进程A} \rightarrow \boxed{\text{打印机}} \leftarrow \text{进程B}
    \]
    一方占用，另一方阻塞等待。
    \vspace{0.5em}
    SPOOLing 方式：
    \[
    \text{进程A} \rightarrow \boxed{\text{输出井}} \searrow
    \]
    \[
    \text{进程B} \rightarrow \boxed{\text{输出井}} \rightarrow
    \boxed{\text{守护进程}} \rightarrow \boxed{\text{打印机}}
    \]
    \[
    \text{进程C} \rightarrow \boxed{\text{输出井}} \nearrow
    \]
    \vspace{0.5em}
    多个进程可同时向输出井写数据，互不阻塞。
\end{frame}

\begin{frame}{SPOOLing 的特点}
    \begin{center}
        \begin{tabular}{|l|l|}
            \hline
            \textbf{特性} & \textbf{说明} \\
            \hline
            设备共享 & 将独占设备改造为共享设备 \\
            虚拟化 & 一台物理设备对应多个虚拟设备 \\
            并行性 & 多进程 I/O 可并行进行 \\
            速度匹配 & 磁盘高速缓冲缓和速度差异 \\
            用户透明 & 用户感知不到 SPOOLing 的存在 \\
            \hline
        \end{tabular}
    \end{center}
    \vspace{0.5em}
    核心价值：
    \[
    \boxed{\text{物理独占} \rightarrow \text{逻辑共享}}
    \]
\end{frame}

\begin{frame}{虚拟设备技术}
    \textbf{虚拟设备}：通过 SPOOLing 技术，一台物理 I/O 设备可对应多个虚拟设备。
    \vspace{0.5em}
    \begin{itemize}
        \item 每个用户进程认为自己独占一台设备
        \item 实际上多进程共享同一台物理设备
        \item OS 负责调度，确保互不干扰
    \end{itemize}
    \vspace{0.5em}
    \begin{center}
        \begin{tabular}{|l|l|l|}
            \hline
            & \textbf{独占设备} & \textbf{共享设备} \\
            \hline
            同时使用 & 仅一个进程 & 多个进程 \\
            典型代表 & 打印机、磁带机 & 磁盘、SSD \\
            分配方式 & 静态分配 & 动态分配 \\
            SPOOLing 后 & $\rightarrow$ 虚拟共享 & --- \\
            \hline
        \end{tabular}
    \end{center}
    \vspace{0.5em}
    SPOOLing 是虚拟设备技术的典型实现。
\end{frame}

\begin{frame}{SPOOLing 使用场景}
    SPOOLing 适用于\textbf{低速独占设备}：
    \vspace{0.5em}
    \begin{description}
        \item[打印机] 最常见的 SPOOLing 应用
        \item[终端输入] 预先将输入数据读入输入井
        \item[网络打印] 远程打印服务
        \item[批处理系统] 早期 SPOOLing 起源于批处理
    \end{description}
    \vspace{0.5em}
    \textbf{不适用}的场景：
    \begin{itemize}
        \item 高速设备（磁盘本身已是共享设备）
        \item 实时设备（需立即响应，不能缓冲延迟）
    \end{itemize}
\end{frame}

\begin{frame}{本节总结}
    \begin{enumerate}
        \item SPOOLing = 假脱机操作，用磁盘模拟独占设备
        \item 系统组成：输入/输出井（磁盘）+ 输入/输出进程 + 缓冲区
        \item 打印机 SPOOLing：输出井暂存 $\rightarrow$ 守护进程实际打印
        \item 将独占设备改造为共享设备（虚拟设备技术）
    \end{enumerate}
    \vspace{1em}
    \begin{center}
        \textcolor{blue}{\textbf{恭喜！操作系统模块全部学习完毕！}}
    \end{center}
    \begin{center}
        \textcolor{red!80!black}{\textbf{下一阶段：进入习题训练与真题实战！}}
    \end{center}
\end{frame}

\end{document}
